\subsection{The same parameter on the eight-point branch quotient}

\begin{result}{E.189}{Proved}
{Four pair comparisons and the common residual branch permutation}
Use the single corrected twelve-point action of E.188.
Write its eight quotient labels as
\[
 B=\{b_0,b_1,b_2,b_3\},\qquad
 J=\{t_1,t_2,t_3,t_4\},
\tag{E.189.1}
\]
where \(b_0\) is the original origin and \(t_j\) is the
image of the \(j\)-th named opposite three-torsion pair.
For the original normalized \(F=(1,f)\) of E.181, the
induced outer automorphism of the actual group
\(\widehat\Gamma_{0,8}\) is
\[
                            (F,\sigma_2)
       \in\widehat{\mathrm{GT}}\times\mathfrak S_8.
\tag{E.189.2}
\]
Here \(\sigma_2\) is the same permutation of the three
new branch labels from E.188; it fixes \(b_0\) and
all four \(t_j\).
The four actual forgetting maps retaining \(B\cup\{t_j\}\)
have the original parameter \(F\).  Forgetting instead
\(b_1,b_2,b_3\) gives exactly the original outer action
with parameter \(F\) on
\(\widehat\Gamma_{0,\{b_0,t_1,t_2,t_3,t_4\}}\).
These statements use one completed marking and its
images under the stated maps.
They do not assert an original finite compiler bound
or yet perform the final CM character calculation.
\end{result}

\begin{proof}
E.188 has established the actual full completed group
on which the action is defined.  The classification used
at this point must not omit a symmetry hypothesis from
the original Harbater--Schneps theorem.
We use the stronger published result of
\citet[Corollary C, p.4; Definition 2.7 and
Corollary 2.8 with proof, pp.31--32 in the
author version]{LitHoshiMinamideMochizuki2022}.
With the set of primes equal to all primes and with five
ordered variable points on the triply punctured sphere,
it states
\[
       \operatorname{Out}(\widehat\Gamma_{0,8})
                    =\widehat{\mathrm{GT}}\times\mathfrak S_8
\tag{E.189.3}
\]
without auxiliary inertia or symmetry assumptions.
Its permutation factor is the actual geometric
permutation of the eight labels.
Its GT factor is identified naturally by forgetting to
five marks.  The identification with the standard
parameter action is the one fixed by
\citet[Main Theorem; \S3.1, Lemma 1 and Theorem 5,
author pp.30--31]{SymHarbaterSchneps2000};
in particular it uses the same positive braid convention
as E.183.

Write the induced automorphism as \((F',\sigma)\).
There is a direct way to determine the entire permutation,
which also retains the distinguished origin.
The action extends to the colored mixed sphere group
\(\widehat D\) in E.188.1, and its action on the quotient
\(\mathfrak S_B\times\mathfrak S_J\) is conjugation by
\(\sigma_2\), by the actual twelve-label calculation.
Conjugation of this geometric permutation subgroup in
the outer automorphism group of the pure subgroup
therefore gives
\[
 \sigma\tau\sigma^{-1}
       =\sigma_2\tau\sigma_2^{-1}
           \qquad(\tau\in\mathfrak S_B\times\mathfrak S_J).
\tag{E.189.4}
\]
The GT factor disappears here because E.189.3 is a
direct product.  The embedding of the geometric
permutation subgroup is faithful for eight points,
as is also part of the stated reconstruction.

The centralizer of
\(\mathfrak S_B\times\mathfrak S_J\) in \(\mathfrak S_8\)
is trivial.  Indeed an element commuting with
\(\mathfrak S_B\) preserves its set \(J\) of fixed
points, and similarly preserves \(B\).  On each
four-element set it centralizes the full symmetric
group, whose center in that natural permutation
representation is trivial.  Equation E.189.4 thus gives
\[
                               \sigma=\sigma_2.
\tag{E.189.5}
\]
This conclusion is stronger than merely preserving
the branch quartet: the original origin is fixed,
and only the three newly added fixed labels can move.

For each \(j\), consider the literal map
\[
 p_j:\widehat\Gamma_{0,8}
                 \longrightarrow
       \widehat\Gamma_{0,B\cup\{t_j\}}.
\tag{E.189.6}
\]
In the surface cover it is obtained by retaining the
selected opposite pair of torus marks and forgetting
the other torus marks.  The four fixed points of the
resulting elliptic involution supply its four branch
labels, and ordering them gives exactly the
five-point subgroup in E.183.
This is an equality of the ordinary quotient and
forgetting maps before completion: a lift of a colored
sphere homeomorphism descends to that same homeomorphism
after the unwanted pair marks are removed.
The covering character is still one at each of the
four branch points, so no unmarked replacement of the
double cover occurs in this square.

For clarity, the mixed two-point torus group causes no
central ambiguity in this comparison.  Its involution
exchanges the two selected poles, has order two, and
commutes with the three lifted braid generators of
E.183.  Those generators generate the ordered subgroup
\(\Gamma_{1,2}=B_4/\langle\Omega\rangle\).
Every mixed class is an ordered class or that class
times the involution.  Consequently
\[
 \Gamma_{1,[2]}=\Gamma_{1,2}\times\langle\iota_2\rangle,
 \qquad
 \Gamma_{1,[2]}/\langle\iota_2\rangle=\Gamma_{1,2},
\tag{E.189.7}
\]
with the displayed ordered-pole marking.
All groups here retain their full finite-index
topologies.  The branch-ordered subgroup is the
same normal index-\(24\) subgroup
\(\Gamma_{0,5}=P_4/\langle\Omega\rangle\).

The corrections of E.188 were made once at twelve
points, and preserve each original selected pole.
Their images after this ordinary forgetting therefore
give an inner correction of the ordered two-point
action of E.182.  By E.183 that action has precisely
the original parameter \(F\) on its branch-ordered
five-point subgroup.  An inner change in the colored
mixed sphere group can contribute a geometric
permutation of the four branch labels; it cannot
alter its GT factor in the direct product.
Thus each of the four maps E.189.6 has induced GT
factor \(F\).

On the other hand, the same factor for \((F',\sigma_2)\)
is \(F'\), by the natural forgetting isomorphism
in the reconstruction theorem (or the standard
forgetting compatibility of
\cite[\S3.1, Lemma 1(c), author p.30]{SymHarbaterSchneps2000}).
Since the GT action on the five-point pure group is
faithful, \(F'=F\).  This proves E.189.2, and it also
shows that all four uses of the two-point comparison
are consistent with the single eight-point action.

Finally forget \(b_1,b_2,b_3\).
The standard GT factor commutes with this ordinary
forgetting map.  The permutation \(\sigma_2\) may be
represented by a braid in a disc containing just
those three marks.  After they are forgotten, that
representative is isotopic to the identity.  A
different representative changes only the final
inner automorphism.  The induced outer action on
the retained five-point group therefore has exactly
the original parameter \(F\), with \(b_0\) retained.
No independently chosen conjugator at one of the
four intermediate targets is used in this argument.
\end{proof}
